% SPDX-License-Identifier: CC-BY-4.0
% Copyright 2025 Florian Aram Feuerriegel — kassensturz.org
% Abschnitt 1 — Einleitung

\section{Einleitung}
\label{sec:einleitung}

Finanzpolitische Entscheidungen verschieben mehrere Größen zugleich: den
Haushaltssaldo, die Verteilung der verfügbaren Einkommen, die Emissionen und
die Wirtschaftsleistung. Ein höherer Rentenbeitrag verbessert den Saldo und
belastet die unteren Einkommen, ein höherer Grundfreibetrag entlastet fast
alle Haushalte und vergrößert das Defizit, ein höherer \COz-Preis senkt die
Emissionen und wirkt regressiv. Im Planspiel \emph{Kassensturz} treten diese
Zielkonflikte gleichzeitig auf: Teams aus vier Ressorts entscheiden über
Steuer-, Sozial- und Klimapolitik und tragen die Folgen über fünf
Legislaturperioden.

Das Modell dahinter muss fünf Anforderungen erfüllen. Es bildet die
Instrumente der deutschen Finanzpolitik auf Einnahmen- und Ausgabenseite ab,
verbindet Haushalte und Staat konsistent, schreibt mehrere Perioden fort,
rechnet schnell genug, um auf jede Entscheidung sofort zu antworten, und nennt
für jede Annahme eine Quelle. Die etablierten
Steuer-Transfer-Mikrosimulationsmodelle für Deutschland -- STSM
\parencite{steiner2012}, IZA$\Psi$MOD \parencite{peichl2010}, das Modell des
ifo Instituts \parencite{bloemer2020} und EUROMOD \parencite{sutherland2013}
-- rechnen mit Tausenden Haushalten aus Befragungsdaten wie SOEP und
EU-SILC, die meisten mit strukturell geschätztem Arbeitsangebot. Sie sind das
richtige Werkzeug für die Politikbewertung
\parencite{bourguignon2006}, aber für ein Spiel weder schnell noch durchschaubar
genug. Makro-Mikro-Verknüpfungen \parencite{bourguignon2008} lösen die
Rückwirkung der Gesamtwirtschaft auf die Haushalte, nicht die Rechenzeit.

Kassensturz geht den umgekehrten Weg: zwölf repräsentative Haushaltstypen,
Verhalten in reduzierter Form über Elastizitäten aus der Literatur und ein
einfacher Makrorahmen. Drei Eigenschaften unterscheiden das Modell von einer
bloßen Vereinfachung. Erstens bucht der Staat, was bei den Haushalten
ankommt: Transfers, Renten, \COz-Last und die Inzidenz der Unternehmen- und
Vermögensteuern summieren sich über die Haushaltstypen zur Buchung des
Staates, und Tests erzwingen diese Identität. Zweitens ist der Status quo der
Rechtsstand 2026 mit seinen beschlossenen Pfaden: Senkung der
Körperschaftsteuer, steigender Rentenbeitrag, Verteidigungsausgaben,
Sondervermögen Infrastruktur und die Schuldenbremse in der Fassung von 2025.
Drittens sind die Wirkungen der Stellgrößen gemessen, nicht behauptet: Jede
Stellgröße wird um einen festen Schritt verschoben, das Modell rechnet den
ganzen Pfad neu, und jede Saldowirkung zerfällt in mechanische Wirkung,
Verhaltensreaktion und Nachfrage. Alle Tabellen, Abbildungen und Kennzahlen der
Abschnitte~\ref{sec:kalibrierung} und~\ref{sec:wirkungen} entstehen aus dem
Modellcode.

Bezogen auf eine Milliarde Saldoverbesserung ordnen sich die Stellgrößen
entlang zweier Preise. Einen Verteilungspreis haben Beitragssätze, Umsatzsteuer,
\COz-Preis und Kürzungen bei Renten und Transfers; Steuern auf hohe
Einkommen, Kapital und Vermögen senken die Ungleichheit. Einen Wachstumspreis
haben alle, am geringsten die Steuern auf Haushalte mit niedriger
Konsumneigung, langfristig am höchsten die Körperschaftsteuer, weil sie über
das private Kapital wirkt.
